%=============================================================================!
% 13.6  RESCALING THAT SAMPLES                                                !
%=============================================================================!
\section{Rescaling that samples}
\label{sec:th-csvr}

Berendsen coupling disturbs the motion little but narrows the spread of
$K$; Langevin's friction gives the right spread but disturbs every atom.
Bussi, Donadio and Parrinello found a rescaling of all the velocities
together, like Berendsen's, that samples the canonical distribution
exactly\cite{Bussi2007}: canonical sampling through velocity rescaling,
CSVR. This section derives it from the exact step of
\cref{sec:th-langevin}, measures it, and reads TrajCast's implementation,
which holds its learned runs at a temperature.

\subsection*{The derivation}

In the $N_{\mathrm f}$ free components of \cref{sec:sm-equipartition},
mass-weighted so that $K = \frac12\lvert\vb{u}\rvert^2$, the velocities
form a vector $\vb{u}$. With the noise of \cref{eq:th-fdr}, the exact step
\cref{eq:th-ou} takes a velocity component $v$ to $cv + \sqrt{(1 -
c^2)\kB T_0/m}\,g$, and so $u = \sqrt m\,v$ to $cu + \sqrt{(1 - c^2)\kB
T_0}\,g$, whatever the mass. Applied to every component with the friction
$\gamma = 1/(2\tau_{\mathrm T})$, chosen so that the mean of $K$ relaxes
with the time constant $\tau_{\mathrm T}$, as shown below, and for which
$c^2 = \mathrm{e}^{-\dt/\tau_{\mathrm T}} = c_1$, it gives
\begin{equation*}
   \vb{u}' = \sqrt{c_1}\,\vb{u} + \sqrt{(1 - c_1)\kB T_0}\;\vb{g} ,
\end{equation*}
with $\vb{g}$ a vector of $N_{\mathrm f}$ independent normal numbers. Its
squared length is
\begin{equation*}
   \lvert\vb{u}'\rvert^2 = c_1\lvert\vb{u}\rvert^2 + (1 - c_1)\kB T_0\,
   \lvert\vb{g}\rvert^2 + 2\sqrt{c_1(1 - c_1)\kB T_0}\;\vb{u}\cdot\vb{g} .
\end{equation*}
The vector $\vb{g}$ can be resolved along $\vb{u}$ and at right angles to
it. Its density, the product of the factors $\mathrm{e}^{-g_k^2/2}$, is a
constant times $\mathrm{e}^{-\lvert\vb{g}\rvert^2/2}$, which depends only on
the length of $\vb{g}$; turning the axes, so that the first lies along
$\vb{u}$ and the others at right angles to it and to one another, leaves
that length, and so the density, unchanged, and the components of
$\vb{g}$ along the new axes are again independent normal numbers. So
$\vb{u}\cdot\vb{g} = \lvert\vb{u}\rvert R_1$, with $R_1$ the component of
$\vb{g}$ along $\vb{u}$, and $\lvert\vb{g}\rvert^2 = R_1^2 + R_\perp^2$,
with $R_\perp^2$ the sum of the squares of the other $N_{\mathrm f} - 1$ components,
independent of $R_1$. Dividing by $\lvert\vb{u}\rvert^2 = 2K$, which turns
the last term into $2R_1\sqrt{c_1(1 - c_1)\kB T_0/2K}$, and writing $c_2
= (1 - c_1)\kB T_0/2K = (1 - c_1)K_0/(N_{\mathrm f}K)$ (the names of Bussi
and colleagues, not the components $c_k$ of \cref{sec:sm-equipartition}),
\begin{equation}
   \frac{K'}{K} = \lambda^2 = c_1 + c_2\bigl(R_1^2 + R_\perp^2\bigr)
   + 2R_1\sqrt{c_1c_2} .
   \label{eq:th-csvr}
\end{equation}
CSVR keeps the change of $K$ that this step makes and throws away the
change of direction: it multiplies every velocity by $\lambda$. Since
$\lambda^2$ depends on the velocities only through $K$, the new $K$ has,
for a given $K$, the same distribution as after the whole
Ornstein-Uhlenbeck step. That step keeps the canonical distribution, every
component Gaussian with variance $\kB T_0$, and with it the gamma
distribution of $K$; so CSVR keeps that distribution too, for any $\dt$
and any $\tau_{\mathrm T}$. In the canonical density the direction of
$\vb{u}$ is spread evenly, independent of its length and of the
positions, and the scaling changes only the length; so, as with
Andersen's redraws, the whole canonical density is kept, and a run whose
motion is ergodic samples it to the accuracy of the step that moves the
atoms. Averaged over the noise
(\cref{ex:th-csvr-mean}), $\ens{K'} - K = (1 - c_1)(K_0 - K)$, the change
that $\dd K/\dd t = (K_0 - K)/\tau_{\mathrm T}$ makes over $\dt$; for a
short step $1 - c_1 \approx \dt/\tau_{\mathrm T}$, and this is the
relaxation of Berendsen coupling with the same time constant. The noise
supplies the canonical spread that Berendsen lacks. The sum $R_\perp^2$ of
$N_{\mathrm f} - 1$ squares need not be drawn square by square: by
\cref{eq:sm-gamma} with $\kB T = 1$ (\cref{ex:th-chi}), half of it has a
density proportional to $y^{(N_{\mathrm f} - 1)/2 - 1}\mathrm{e}^{-y}$,
the gamma distribution of shape $(N_{\mathrm f} - 1)/2$, from which NumPy
draws directly. \texttt{CSVR.before} is \cref{eq:th-csvr}:

\begin{code}
        k = kinetic(m, v, self.mv2_to_energy)
        target = 0.5 * self.n_free * self._kt(self.temperature)
        c1 = math.exp(-dt / self.tau)
        c2 = (1 - c1) * target / (k * self.n_free)
        r1 = rng.standard_normal(k.shape)
        rest = (2 * rng.standard_gamma(0.5 * (self.n_free - 1), k.shape)
                if self.n_free > 1 else np.zeros(k.shape))
        alpha2 = c1 + c2 * (rest + r1 * r1) + 2 * r1 * np.sqrt(c1 * c2)
        return v * np.sqrt(alpha2)[..., None, None]
\end{code}

\noindent Bussi and colleagues wrote the factor as $\alpha$, which is
$\lambda$ in this chapter; \texttt{alpha2} holds its square. At \SI{150}{\kelvin},
with a target of \SI{135}{\kelvin}, $\tau_{\mathrm T} =
\SI{100}{\femto\second}$ and steps of \SI{10}{\femto\second}, $c_1 =
0.90484$ and $c_2 = \num{1.120e-4}$, and with $R_1 = 0$ and $R_\perp^2$ at its
mean of 764, $\lambda^2 = 0.99037$, close to Berendsen's 0.9900.

\subsection*{Measured}

Over three runs of liquid argon the spread of $K$ is 0.0499 with
$\tau_{\mathrm T} = \SI{100}{\femto\second}$ and 0.0508 with
\SI{1}{\pico\second}, against the canonical 0.0511
(\cref{fig:th-berendsen}(a)); the diffusion coefficient is 1.07 and 1.05
times its value at fixed energy, a gap that \cref{sec:th-compare} weighs
against the scatter of the runs.
Acting once every three steps, over their whole length, with
$\tau_{\mathrm T} = \SI{300}{\femto\second}$, it still gives 0.0513.
For a single lithium atom, the density of the potential energy is within
0.012 of the exact one, and the atoms hop 1.08 times per picosecond. An
isolated molecule would show another side of the common factor: it keeps
its angular momentum under the forces, and CSVR's scaling multiplies it
at random, so TrajCast's authors set it to zero every hundred strides for
paracetamol\cite{Thiemann2026}.
CSVR scales every velocity by one factor, as Berendsen coupling does, but
its factor is not tilted towards growth\cite{Braun2018}, and a motion of
the centre of mass left in the velocities does not grow
(\cref{sec:th-trajectory}).

\subsection*{TrajCast's thermostat}

TrajCast predicts the state a stride $\Dt$ ahead, from
\SI{5}{\femto\second} for water to \SI{30}{\femto\second} for quartz in
the runs reported, and after each stride applies CSVR once, in the
file \texttt{\_thermostat.py} of its forecasting tools\cite{Thiemann2026},
whose lines 96-109 are \cref{eq:th-csvr}:

\begin{code}
        c2 = (
            (torch.tensor(1.0) - self.c1) * e_kin_target
            / e_kin_current / self.n_dofs
        ).to(self.device)
        r1 = torch.randn(1, device=self.device)
        r2 = self._draw_sum_noises_from_gamma_dist(n_dofs_1=self.n_dofs - 1)
        alpha_2 = (
            self.c1
            + c2 * (r1 * r1 + r2)
            + torch.tensor(2.0) * r1 * (self.c1 * c2).sqrt()
        )
\end{code}

\noindent with comments removed and one line wrapped. \texttt{c1} is $\mathrm{e}^{-\Dt/\tau_{
\mathrm T}}$, computed once with the stride as the step (lines 32-37);
\texttt{r2} is $R_\perp^2$, drawn as twice a gamma number of shape $(N_{\mathrm f} -
1)/2$, or, when $N_{\mathrm f} - 1$ is odd, as twice one of shape
$(N_{\mathrm f} - 2)/2$ plus one more square (lines 113-123), which is the
same distribution; and \texttt{n\_dofs} is $3N$ less three for each kind
of momentum, linear or angular, that the run removes or holds fixed, or
less a number the run's settings give (\texttt{forecast.py}, lines
344-378).
With no time constant given, $\tau_{\mathrm T}$ is a hundred strides, so
$c_1 = \mathrm{e}^{-0.01} = 0.9900$ and the temperature is held loosely;
the runs reported used ten strides for water and paracetamol, 50 and
\SI{70}{\femto\second}, $c_1 = 0.9048$, and two and a half for quartz,
\SI{75}{\femto\second}, $c_1 = 0.6703$. Since \cref{eq:th-csvr} keeps the
distribution of $K$ for any step, the thermostat itself works at any
stride, as the run acting once every three steps confirmed. The stride
changes what the thermostat corrects: a learned step may not keep the
energy, and CSVR would then remove the error as heat, so in a learned run
the heat may measure the model's error as well as the bath's work, a
theme of the planned Part~VI. TrajCast's authors found that the longer strides needed
much shorter $\tau_{\mathrm T}$ than usual to hold the temperature and its
spread, and put it down to each stride standing for several steps taken
before one rescaling\cite{Thiemann2026}.

\begin{keyidea}
CSVR applies the exact Ornstein-Uhlenbeck step to all $N_{\mathrm f}$
mass-weighted velocity components at once, keeps only its change of
length, and scales every velocity by the factor $\lambda$ of
\cref{eq:th-csvr}. Its step then keeps the canonical gamma distribution of
$K$ for any step and coupling, the mean relaxes as under Berendsen
coupling, and the motion is disturbed as little as by Berendsen
coupling. TrajCast applies it once per
learned stride, with $c_1 = \mathrm{e}^{-\Dt/\tau_{\mathrm T}}$.
\end{keyidea}

\notebook{13}{set the coupling and the stride, and compare CSVR with Berendsen}

\begin{selfcheck}
For TrajCast's stride of \SI{30}{\femto\second} for quartz, what
$\tau_{\mathrm T}$ gives $c_1 = 0.5$?
\end{selfcheck}
